\documentclass[final]{beamer}

\usepackage[size=a0,orientation=landscape,scale=1.16]{beamerposter}
\usepackage[T1]{fontenc}
\usepackage[utf8]{inputenc}
\usepackage{lmodern}
\usepackage{booktabs}
\usepackage{array}
\usepackage{tikz}
\usepackage{ragged2e}

\definecolor{PosterInk}{HTML}{172019}
\definecolor{PosterMuted}{HTML}{5B645D}
\definecolor{PosterPaper}{HTML}{F7F0DF}
\definecolor{PosterCream}{HTML}{FFFAF0}
\definecolor{PosterGreen}{HTML}{0F7F64}
\definecolor{PosterGreenSoft}{HTML}{D8EFE3}
\definecolor{PosterClay}{HTML}{C8643D}
\definecolor{PosterClaySoft}{HTML}{F4D5C6}
\definecolor{PosterBlue}{HTML}{315F85}
\definecolor{PosterBlueSoft}{HTML}{D7E5EE}
\definecolor{PosterGold}{HTML}{D4A72C}
\definecolor{PosterGoldSoft}{HTML}{F5E6B5}
\definecolor{PosterCharcoal}{HTML}{18221D}

\setbeamercolor{background canvas}{bg=PosterPaper}
\setbeamercolor{normal text}{fg=PosterInk}
\setbeamercolor{block title}{fg=PosterCream,bg=PosterCharcoal}
\setbeamercolor{block body}{fg=PosterInk,bg=PosterCream}
\setbeamerfont{block title}{size=\Large,series=\bfseries}
\setbeamerfont{block body}{size=\normalsize}
\setbeamertemplate{navigation symbols}{}
\setbeamertemplate{headline}{}
\setbeamertemplate{footline}{}

\newcommand{\softblock}[1]{%
  \setbeamercolor{block title}{fg=PosterCream,bg=PosterCharcoal}%
  \setbeamercolor{block body}{fg=PosterInk,bg=PosterCream}%
  #1%
}

\newcommand{\darkblock}[1]{%
  \setbeamercolor{block title}{fg=PosterCream,bg=PosterCharcoal}%
  \setbeamercolor{block body}{fg=PosterCream,bg=PosterCharcoal}%
  #1%
  \setbeamercolor{block body}{fg=PosterInk,bg=PosterCream}%
}

\newcommand{\metricbar}[4]{%
  \begin{tikzpicture}[baseline=-0.5ex]
    \node[anchor=west, font=\bfseries] at (0,0) {#1};
    \fill[PosterInk!10, rounded corners=4pt] (4.3,-0.18) rectangle (13.2,0.18);
    \fill[#3, rounded corners=4pt] (4.3,-0.18) rectangle ({4.3 + 8.9*#2/100},0.18);
    \node[anchor=east, font=\bfseries] at (15.1,0) {#4};
  \end{tikzpicture}%
}

\newcommand{\smalltag}[1]{%
  \tikz[baseline=(X.base)]\node[
    draw=PosterInk!18,
    fill=PosterCream,
    rounded corners=8pt,
    inner xsep=8pt,
    inner ysep=4pt,
    font=\bfseries\small
  ] (X) {#1};%
}

\title{\vspace{-1.3cm}TA-MPQ: Structured Mixed Precision at an Exact INT4 Budget}
\author{Current snapshot: 2026-04-22}
\institute{Qwen3.5-9B \quad | \quad 2/4/8-bit mixed precision \quad | \quad MMLU-coding selected, EvalPlus validated}

\begin{document}
\begin{frame}[t]

\begin{tikzpicture}[remember picture,overlay]
  \fill[PosterGold!22] (current page.north west) circle [radius=18cm];
  \fill[PosterBlue!16] ([xshift=-4cm,yshift=-2cm]current page.north east) circle [radius=20cm];
  \draw[PosterInk!8, line width=2pt, rounded corners=22pt]
    ([xshift=1.2cm,yshift=-1.2cm]current page.north west)
    rectangle
    ([xshift=-1.2cm,yshift=1.2cm]current page.south east);
\end{tikzpicture}

\vspace{0.6cm}
\begin{columns}[t,totalwidth=\textwidth]
  \begin{column}{0.70\textwidth}
    {\usebeamerfont{title}\color{PosterCharcoal}\inserttitle\par}
    \vspace{0.35cm}
    {\Large\color{PosterMuted}
    We moved from evolutionary crossover to a deterministic task-sensitivity frontier:
    promote the most valuable groups to 8-bit, demote lower-value groups to 2-bit,
    and keep the full model near the raw uniform-INT4 weight footprint.\par}
    \vspace{0.45cm}
    \smalltag{Qwen3.5-9B}\hspace{0.25cm}
    \smalltag{2/4/8-bit policy}\hspace{0.25cm}
    \smalltag{Exact INT4 budget}\hspace{0.25cm}
    \smalltag{BigCodeBench generation running}
  \end{column}
  \begin{column}{0.26\textwidth}
    \begin{beamercolorbox}[rounded=true,shadow=true,wd=\textwidth,sep=0.55cm]{block title}
      {\Large Current mixed policy}\\[0.2cm]
      {\fontsize{54}{58}\selectfont\bfseries 3.999-bit}\\[0.2cm]
      {\large Estimated average bit-width at 3.694 GB with 0.026\% budget slack.}\\[0.25cm]
      {\normalsize \texttt{gpf\_refine\_00\_i0112}\\
      hash: \texttt{7caae6e...2043e4}}
    \end{beamercolorbox}
  \end{column}
\end{columns}

\vspace{0.7cm}

\begin{columns}[t,totalwidth=\textwidth]
  \begin{column}{0.32\textwidth}
    \begin{block}{Why the route changed}
      \justifying
      The old evolutionary route mixed random mutation, crossover, surrogate ranking,
      and budget repair. In practice, crossover fought the exact-budget constraint
      and produced unstable candidates.

      \vspace{0.35cm}
      The new route makes the budget explicit: traverse an ordered frontier from
      mostly INT4 toward more INT8, using INT2 demotions only where needed to stay
      near the target size.

      \vspace{0.4cm}
      \begin{tabular}{@{}p{0.18\linewidth}p{0.76\linewidth}@{}}
      \textbf{1. Profile} & Task sensitivity on MMLU-coding dev prompts.\\
      \textbf{2. Rank} & Group value for 8-bit promotions and 2-bit demotions.\\
      \textbf{3. Build} & Greedy frontier under exact INT4 footprint.\\
      \textbf{4. Refine} & Local search around the best coarse region.\\
      \textbf{5. Validate} & Run real quantized artifacts on held-out tasks.\\
      \end{tabular}

      \vspace{0.45cm}
      \begin{beamercolorbox}[rounded=true,wd=\textwidth,sep=0.35cm]{block body}
        \colorbox{PosterGreenSoft}{\parbox{0.94\textwidth}{
        \textbf{Main claim today.}
        The method is now controllable and reproducible. Performance is mixed:
        it improves over INT4 on HumanEval/EvalPlus, but it does not yet beat
        INT4 or INT8 on MMLU-coding.
        }}
      \end{beamercolorbox}
    \end{block}

    \begin{block}{MMLU-coding, last100}
      \small Simple-evals style prompt, max\_new\_tokens=4096. Single-run last100.
      \vspace{0.35cm}

      \metricbar{BF16}{95}{PosterGreen}{95\%}\\[0.22cm]
      \metricbar{Uniform INT8}{94}{PosterBlue}{94\%}\\[0.22cm]
      \metricbar{Uniform INT4}{91}{PosterClay}{91\%}\\[0.22cm]
      \metricbar{Mixed 2/4/8}{91}{PosterGold}{91\%}

      \vspace{0.45cm}
      \begin{tabular}{lrr}
        \toprule
        Model & Capped & Avg latency\\
        \midrule
        BF16 & 0 & 9.14s\\
        INT8 & 2 & 46.21s\\
        INT4 & 5 & 66.88s\\
        Mixed & 3 & 12.29s\\
        \bottomrule
      \end{tabular}
    \end{block}
  \end{column}

  \begin{column}{0.32\textwidth}
    \darkblock{
    \begin{block}{Current mixed policy}
      \justifying
      The selected policy changes 169 of 249 groups relative to uniform INT4 while
      matching the INT4 raw weight budget.

      \vspace{0.45cm}
      \begin{tikzpicture}
        \fill[PosterClay, rounded corners=8pt] (0,0) rectangle (3.12,0.52);
        \fill[PosterGold] (3.12,0) rectangle (10.44,0.52);
        \fill[PosterBlue, rounded corners=8pt] (10.44,0) rectangle (12,0.52);
        \node[anchor=west, text=PosterCream, font=\bfseries] at (0,0.95) {Param mass: 26.00\% 2-bit \quad 61.02\% 4-bit \quad 12.98\% 8-bit};
      \end{tikzpicture}

      \vspace{0.45cm}
      \begin{tabular}{lrr}
        \toprule
        Bit & Groups & Param mass\\
        \midrule
        2-bit & 57 & 26.00\%\\
        4-bit & 80 & 61.02\%\\
        8-bit & 112 & 12.98\%\\
        \bottomrule
      \end{tabular}

      \vspace{0.35cm}
      \begin{tabular}{lr}
        Total groups & 249\\
        Estimated average bit-width & 3.9989\\
        Estimated footprint & 3.694 GB\\
        Budget slack & 0.026\%\\
      \end{tabular}

      \vspace{0.3cm}
      {\small Budget accounting is raw linear weight footprint, not end-to-end VRAM or throughput parity.}
    \end{block}
    }

    \begin{block}{HumanEval / EvalPlus}
      \small Official EvalPlus evaluator. pass@1 with greedy generation, max\_new\_tokens=768.
      \vspace{0.35cm}

      \begin{tabular}{lrr}
        \toprule
        Model & Base pass@1 & HumanEval+ pass@1\\
        \midrule
        BF16 & 87.8\% & 82.3\%\\
        Uniform INT8 & 89.0\% & 82.3\%\\
        Uniform INT4 & 81.1\% & 76.8\%\\
        \textbf{Mixed 2/4/8} & \textbf{84.8\%} & \textbf{79.9\%}\\
        \bottomrule
      \end{tabular}

      \vspace{0.45cm}
      \colorbox{PosterGreenSoft}{\parbox{0.94\textwidth}{
      \textbf{Transfer signal.}
      Mixed recovers +3.7 points over INT4 on HumanEval base and +3.0 points on
      HumanEval+, while keeping the INT4-sized budget.
      }}
    \end{block}
  \end{column}

  \begin{column}{0.32\textwidth}
    \begin{block}{BigCodeBench status}
      \justifying
      We use the official BigCodeBench instruct prompt builder, not simple-evals.
      Generation is separated from execution so evaluator failures do not erase samples.

      \vspace{0.4cm}
      \begin{tabular}{lrrl}
        \toprule
        Model & Generated & Capped & Status\\
        \midrule
        BF16 & 148 / 148 & 2 & Complete; eval pending\\
        Uniform INT8 & 20 / 148 & 0 & H100 running\\
        Uniform INT4 & 20 / 148 & 0 & H100 running\\
        Mixed 2/4/8 & 30 / 148 & 1 & H100 running\\
        \bottomrule
      \end{tabular}

      \vspace{0.35cm}
      {\small BigCodeBench scores are intentionally not reported yet; generation artifacts exist, execution grading is pending.}
    \end{block}

    \darkblock{
    \begin{block}{What we can say now}
      \begin{itemize}
        \item INT8 remains close to BF16 across current coding evaluations.
        \item Uniform INT4 is a strong size baseline, but loses clear accuracy on HumanEval/EvalPlus.
        \item The mixed exact-budget policy improves over INT4 on HumanEval while matching the INT4 footprint.
        \item MMLU-coding currently does not separate mixed from INT4; stronger selection and evaluation coverage is needed.
        \item BigCodeBench is the next meaningful generative-code test once execution grading is stable.
      \end{itemize}
    \end{block}
    }

    \begin{block}{Next steps}
      \begin{itemize}
        \item Finish BigCodeBench generation for INT8, INT4, and mixed; then run official execution grading.
        \item Report end-to-end VRAM and throughput separately from raw weight footprint.
        \item Use BigCodeBench results to decide whether MMLU-coding sensitivity is sufficient or needs task-specific refinement.
        \item If BigCodeBench shows mixed gains, promote this route as the main TA-MPQ story; if not, tune the sensitivity profiler and candidate generator.
      \end{itemize}
    \end{block}
  \end{column}
\end{columns}

\vfill
\begin{center}
{\large\color{PosterMuted}
Artifacts: \texttt{gpf\_refine\_00\_i0112}; BigCodeBench generation-only pipeline with checkpointed samples.}
\end{center}

\end{frame}
\end{document}
